\documentclass[11pt,a4paper]{article}
% Shared preamble for the English probability and statistics documents.
% Based on the original course preamble. Compile with pdfLaTeX.
\usepackage[top=1.5cm,bottom=2cm,left=2cm,right=2cm]{geometry}
\usepackage[T1]{fontenc}
\usepackage{lmodern}
\usepackage[utf8]{inputenc}
\usepackage[english]{babel}

\usepackage{amsmath}
\usepackage{amssymb}
\usepackage{amsthm}
\usepackage{mathrsfs}
\usepackage{mathtools}
\usepackage{etoolbox}
\usepackage{tikz}
\usetikzlibrary{shapes.multipart,arrows.meta}
\usepackage{xcolor}
\usepackage{enumerate}
\usepackage[shortlabels]{enumitem}
\usepackage{hyperref}
\usepackage{version}
\usepackage{bbold}
\usepackage{fontawesome5}
\usepackage{fancyhdr}
\lfoot{}
\usepackage{natbib}
\usepackage{graphicx}
\usepackage{pdfpages}
\usepackage{authblk}
    
\usepackage{blindtext}
\usepackage{multicol}
\usepackage{collcell}
\usepackage{wrapfig}
\usepackage{blkarray}
\usepackage{cancel}
\usepackage{appendix}
\usepackage{float}
\makeatletter
\input{tex/assets/slashbox.sty}
\makeatother
\hypersetup{unicode,colorlinks,citecolor=red,filecolor=black,
  linkcolor=red,urlcolor=cyan}
\setlength{\columnsep}{1cm}
\pagenumbering{gobble}

\def\numtext#1{%
  \ifcase#1\or one\or two\or three\or four\or five\or six\or seven\or eight\or nine\or ten\or
  eleven\or twelve\or thirteen\or fourteen\or fifteen\or sixteen\or seventeen\or eighteen\or nineteen\or twenty\or Lots\fi}

\newcommand{\Byte}[1]{%
  \def\nodecontents{}%
  \foreach[count=\l] \k in {#1} {\xappto{\nodecontents}{\noexpand\nodepart{\numtext{\l}}\k}}%
  \node[rectangle split,rectangle split parts=8,rectangle split horizontal,draw] at (2,2){\nodecontents};}

\newcommand{\D}{{\mathrm{d}}}
\newcommand{\cE}{{\mathcal{E}}}
\newcommand{\cG}{{\mathcal{G}}}
\newcommand{\cX}{{\mathcal{X}}}
\newcommand{\cY}{{\mathcal{Y}}}
\newcommand{\cB}{{\mathcal{B}}}
\newcommand{\cF}{{\mathcal{F}}}
\newcommand{\cO}{{\mathcal{O}}}
\newcommand{\cS}{{\mathcal{S}}}
\newcommand{\cT}{{\mathcal{T}}}
\newcommand{\given}{\, \big| \,}
\newcommand{\cL}{{\mathscr{L}}}
\newcommand{\cD}{{\mathcal{D}}}
\newcommand{\cA}{{\mathscr{A}}}
\newcommand{\cP}{{\mathscr{P}}}
\newcommand{\cR}{{\mathscr{R}}}
\newcommand{\A}{{\mathcal A}}
\newcommand{\one}{\mathbb{1}}
\newcommand{\mathbbm}[1]{\mathbb{#1}}

\providecommand{\C}{}
\renewcommand{\C}{\mathbb C}
\newcommand{\R}{\mathbb R}
\newcommand{\Q}{\mathbb Q}
\newcommand{\N}{\mathbb N}
\newcommand{\T}{\mathbb T}
\newcommand{\F}{\mathbb F}
\renewcommand{\P}{\mathbb P}
\newcommand{\Z}{\mathbb Z}
\newcommand{\E}{\mathbb {E}}
\newcommand{\Pbb}[1]{\P\left( #1\right)}
\newcommand{\Ebb}[1]{\E\left[ #1\right]}

\newcommand{\lb}{\left (}
\newcommand{\rb}{\right )}
\newcommand{\lbb}{\left [}
\newcommand{\rbb}{\right ]}
\newcommand{\labs}{\left |}
\newcommand{\rabs}{\right |}
\newcommand{\lbrb}[1]{\lb #1 \rb}
\newcommand{\lbbrbb}[1]{\lbb#1\rbb}
\newcommand{\labsrabs}[1]{\labs#1\rabs}
\newcommand{\lbbrb}[1]{\lbb#1\rb}
\newcommand{\lbrbb}[1]{\lb#1\rbb}
\newcommand{\langlerangle}[1]{\langle#1\rangle}
\newcommand{\lbcurly}{\{}
\newcommand{\rbcurly}{\}}
\newcommand{\curly}[1]{\lbcurly#1\rbcurly}
\newcommand{\abs}{\labsrabs}
\newcommand{\ind}[1]{1_{#1}}
\newcommand{\indep}{\perp \!\!\! \perp}
\renewcommand{\d}{\, \mathrm{d}}

\newtheorem{theorem}{Theorem}
\newtheorem{problemCnt}{Lemma}
\theoremstyle{definition}
\newtheorem{problem}[problemCnt]{Problem}
\newtheorem{problembook}[problemCnt]{\faBook \,  Problem}
\newtheorem{problempencil}[problemCnt]{\faPencil* \,  Problem}
\newtheorem{problemstar}[problemCnt]{\faStar\,  Problem}
\newtheorem{problemstarpencil}[problemCnt]{\faStar \, \faPencil* \,   Problem}
\newtheorem{probleminteresting}[problemCnt]{\faSmile[regular] \,   Problem}
\newtheorem{solution}{Solution}

\newcommand{\fwcell}[1]{\makebox[\arraystretch\normalbaselineskip]{$#1$}}
\newenvironment{enumeratei}{\begin{enumerate}[(i)]}{\end{enumerate}}
\newenvironment{enumeratea}{\begin{enumerate}[(a)]}{\end{enumerate}}
\newenvironment{enumeraten}{\begin{enumerate}[\upshape 1.]}{\end{enumerate}}

% Compatibility helpers needed by a few files in the combined archive.
\AtBeginDocument{\noindent{\footnotesize English translation prepared with ChatGPT (OpenAI, GPT-6), 26 October 2026, for translation of the Bulgarian text into English and LaTeX formatting.
The Bulgarian originals remain the reference; translation may contain errors.}\par\smallskip}
\newcommand{\MissingGraphic}[1]{%
  \begin{center}\fbox{\parbox{0.82\linewidth}{\centering Missing source image: \texttt{#1}}}\end{center}}



\pagestyle{fancy}\fancyhf{}\lhead{Probability and Statistics: Resit solutions}\rhead{17 August 2023}
\setlength{\headheight}{15pt}\pagenumbering{arabic}\cfoot{\thepage}
\begin{document}
\noindent\emph{English transcription of the original handwritten solution outline. The original level of detail is retained, including the parts described there as standard.}

\section*{Problem 1}
\subsection*{1. Expected winnings}
\begin{align*}
\E\left[\sum_{i=1}^3\left(\mathbf1_{\{\text{coin }i\text{ is heads}\}}
+3\mathbf1_{\{\text{die }i\text{ shows }6\}}\right)\right]
&=3\P(\text{coin 1 is heads})+3\cdot3\P(\text{die 1 shows }6)\\
&=\frac32+\frac96=3.
\end{align*}
\subsection*{2. A sum divisible by 6}
Whatever the sum before round $i$, exactly one die outcome makes the new sum congruent to 0 modulo 6. Hence the number of rounds has distribution $\operatorname{Ge}(1/6)$. Put $p=1/6$; then
\begin{align*}
\E[\text{winnings}]&=\sum_{k=1}^\infty kp(1-p)^{k-1}\\
&=\frac p{1-p}\sum_{k=1}^\infty k(1-p)^k
=\frac p{1-p}\frac{1-p}{p^2}=\frac1p=6.
\end{align*}
This is precisely $\E[\operatorname{Ge}(1/6)]$.

\subsection*{3. Identifying a distribution}
Since $\P(X=k)/\P(X=k-1)=3/k$,
\begin{align*}
\P(X=k)&=\frac3k\P(X=k-1)
=\frac3k\frac3{k-1}\cdots\frac31\P(X=0)\\
&=\frac{3^k}{k!}\P(X=0).
\end{align*}
Because the probabilities sum to 1,
\[
\P(X=0)\underbrace{\left(1+\frac3{1!}+\frac{3^2}{2!}+\cdots\right)}_{e^3}=1.
\]
Thus $\P(X=0)=e^{-3}$ and $\P(X=k)=e^{-3}3^k/k!$. Hence
\[
X\sim\operatorname{Poi}(3),\qquad\E X=DX=3.
\]

\subsection*{4. Visitors at entrance 1}
Let $X$ be the number of visitors using entrance 1, and let $Y\sim\operatorname{Poi}(\lambda)$ be the total number of visitors. Compute the generating function of $X$:
\begin{align*}
\E s^X&=\sum_{k=0}^\infty\E[s^X\mid Y=k]\P(Y=k)\\
&=\sum_{k=0}^\infty
\E\left[s^{\sum_{i=1}^k\mathbf1_{\{\text{visitor }i\text{ chooses entrance 1}\}}}\right]
e^{-\lambda}\frac{\lambda^k}{k!}\\
&=\sum_{k=0}^\infty
\left(\E s^{\mathbf1_{\{\text{visitor 1 chooses entrance 1}\}}}\right)^k
e^{-\lambda}\frac{\lambda^k}{k!}\\
&=\sum_{k=0}^\infty\left(\frac9{10}s^0+\frac1{10}s^1\right)^k
e^{-\lambda}\frac{\lambda^k}{k!}\\
&=e^{-\lambda}e^{(9/10+s/10)\lambda}
=e^{(\lambda/10)(s-1)}.
\end{align*}
The third equality uses independence. For $Z\sim\operatorname{Poi}(\lambda)$,
\[
\E s^Z=\sum_{k=0}^\infty s^k e^{-\lambda}\frac{\lambda^k}{k!}
=e^{-\lambda}e^{\lambda s}=e^{\lambda(s-1)}.
\]
Comparing generating functions gives
\[
X\sim\operatorname{Poi}(\lambda/10),\qquad\E X=DX=\lambda/10.
\]
In fact, as proved in the lectures, binomial thinning of $\operatorname{Poi}(\lambda)$ with probability $p$ has distribution $\operatorname{Poi}(\lambda p)$.

\section*{Problem 2}
\subsection*{1. Time to a winner}
For $n\ge3$,
\begin{align*}
p_n:=\P(\text{a winner exists})
&=\P(\text{one head, all others tails})
+\P(\text{one tail, all others heads})\\
&=2\frac n{2^n}=\frac n{2^{n-1}}.
\end{align*}
Therefore $X\sim\operatorname{Ge}(p_n)$, with
\[
\E X=\frac{2^{n-1}}n=\frac1{p_n},\qquad DX=\frac{1-p_n}{p_n^2}.
\]
\subsection*{2. A fair entry price}
The total payout is
\[
100[(n-1)+(n-2)+\cdots+2]
=100\left(\frac{(n-1)n}2-1\right)=:S.
\]
By symmetry, the expected winnings of player $i$ do not depend on $i$. Thus
\[
\sum_{i=1}^n\E[\text{winnings of player }i]
=S=n\E[\text{winnings of player 1}].
\]
The fair entry price is therefore $S/n$.
\subsection*{3. Total number of rounds}
With $k$ players remaining, part 1 gives an expected $2^{k-1}/k$ rounds until the next winner. Hence the expected total number of rounds is
\[
\frac{2^{n-1}}n+\frac{2^{n-2}}{n-1}+\cdots+\frac{2^2}3
=\sum_{k=3}^n\frac{2^{k-1}}k.
\]

\section*{Problem 3}
Parts 1, 2 and 3 are standard.
\paragraph{4.} Set $Z_i=X_i-Y_i$. By the central limit theorem,
\begin{align*}
\P(X_1+\cdots+X_n>Y_1+\cdots+Y_n)
&=\P(Z_1+\cdots+Z_n>0)\\
&\approx\P\left(N(0,1)>\frac{-n\E Z_1}{\sqrt{nDZ_1}}\right).
\end{align*}
The required mean and variance are
\[
\E Z_1=\E[X-Y]=\E X-\E Y,
\]
using the quantities calculated in part 1, and
\begin{align*}
DZ_1&=\E[(X-Y)^2]-(\E Z_1)^2\\
&=\E X^2+\E Y^2-2\E[XY]-(\E Z_1)^2.
\end{align*}
The second moments come from part 1, the product moment from part 3 (or directly), and the final squared mean from the preceding calculation.
\end{document}
